\par\Needspace{4\baselineskip}
\begin{center}
\hypertarget{native-066-article048}{}\pdfbookmark[2]{48.}{article-048}48.
\end{center}
Thus, to specify the position of a celestial body in space at any instant, we must know the following quantities for an elliptical orbit.

I. The mean longitude at some arbitrarily chosen instant, called the \emph{epoch}. The longitude itself is sometimes also given this name. The epoch is usually chosen at the beginning of a year, namely noon on 1 January in a leap year or noon on the preceding 31 December in a common year.

II. The mean motion over a specified interval of time, for example one mean solar day or $365$, $365\tfrac{1}{4}$, or $36525$ days.

III. The semimajor axis. It could indeed be omitted whenever the body's mass is either known or can be neglected, since the mean motion already determines the semimajor axis (art. 7). For convenience, however, both are customarily always given.

IV. The eccentricity.

V. The longitude of the perihelion.

\SourceLeaf{067}VI. The longitude of the ascending node.

VII. The inclination of the orbit.

These seven quantities are called the \emph{elements} of the body's motion.

For a parabola or a hyperbola, the time of passage through perihelion takes the place of the first element. For II, we give the quantities analogous to the daily mean motion for these kinds of conic sections (see art. 19 and, for hyperbolic motion, the quantity $\lambda k b^{-\frac{3}{2}}$ in art. 23). For a hyperbola, the remaining elements. can likewise be retained. For a parabola, where the major axis is infinite and the eccentricity $=1$, only the distance at perihelion is given in place of elements III and IV.

\par\Needspace{4\baselineskip}
\begin{center}
\hypertarget{native-067-article049}{}\pdfbookmark[2]{49.}{article-049}49.
\end{center}
In ordinary usage the inclination of the orbit, which we measure from $0$ to $180^0$, extends only to $90^0$. If the angle between the orbit and the arc $\AscendingNode B$ (Fig. 1) exceeds a right angle, the angle between the orbit and the arc $\AscendingNode A$ (which adds to the former angle to give $180^0$) is taken as the inclination of the orbit. In that case, it must be added that the motion is \emph{retrograde} (as when $E\AscendingNode F$ represents part of the orbit in our figure), to distinguish it from the other case, in which the motion is called \emph{direct}. Longitude in the orbit is then usually measured so that at $\AscendingNode$ it agrees with the longitude of this point on the ecliptic, and \emph{decreases} in the direction $\AscendingNode F$. Thus the starting point from which longitudes are measured opposite to the motion lies as far along $\AscendingNode F$ from $\AscendingNode$ as the vernal equinox lies from the same $\AscendingNode$ along $\AscendingNode A$. Hence in this case the longitude of the perihelion is the longitude of the node minus the distance of the perihelion from the node\raisebox{.4ex}{.} Either convention can thus be readily converted into the other. We have preferred ours so that we can dispense with the distinction between direct and retrograde motion and always use the same formulas for both, since ordinary usage very often requires two sets of rules.
